\documentclass[11pt]{article}

%==================== PACKAGES ====================
\usepackage[a4paper,margin=2.2cm,top=2.4cm,bottom=2.4cm]{geometry}
\usepackage[T1]{fontenc}
\usepackage[utf8]{inputenc}
\usepackage{lmodern}
\usepackage{microtype}
\usepackage{xcolor}
\usepackage{titlesec}
\usepackage{enumitem}
\usepackage{fancyhdr}
\usepackage{tcolorbox}
\usepackage{booktabs}
\usepackage{array}
\usepackage{hyperref}
\tcbuselibrary{skins,breakable}

%==================== COLORS ====================
\definecolor{poliblue}{RGB}{0,40,85}    % PoliMi deep blue
\definecolor{polilight}{RGB}{210,225,240}
\definecolor{accent}{RGB}{0,90,160}
\definecolor{codebg}{RGB}{244,246,249}
\definecolor{codeframe}{RGB}{200,210,222}
\definecolor{soft}{RGB}{90,100,110}

%==================== HYPERREF ====================
\hypersetup{
  colorlinks=true,
  linkcolor=accent,
  urlcolor=accent,
  pdftitle={PoliMi Thesis Template -- Quick Guide},
  pdfauthor={Claude Code}
}

%==================== HEADER/FOOTER ====================
\pagestyle{fancy}
\fancyhf{}
\renewcommand{\headrulewidth}{0pt}
\fancyfoot[L]{\small\color{soft}PoliMi MSc Thesis Template}
\fancyfoot[R]{\small\color{soft}\thepage}

%==================== SECTION STYLES ====================
\titleformat{\section}
  {\Large\bfseries\color{poliblue}}
  {}{0pt}
  {\colorbox{poliblue}{\color{white}\,\thesection\,}\hspace{8pt}}
  [\vspace{2pt}{\color{polilight}\titlerule[2pt]}]
\titleformat{\subsection}
  {\large\bfseries\color{accent}}{\thesubsection}{10pt}{}
\titlespacing*{\section}{0pt}{16pt}{8pt}

%==================== CUSTOM COMMANDS ====================
% callout marker (replaces fontawesome icons)
\newcommand{\marker}[2]{\textcolor{#1}{\rule[0.05em]{0.7em}{0.7em}}~\textbf{\textcolor{#1}{#2}}}

% inline code
\newcommand{\code}[1]{%
  \colorbox{codebg}{\rule[-0.2em]{0pt}{1em}\texttt{\small #1}}}

% code block
\newtcolorbox{cmd}{
  enhanced, breakable,
  colback=codebg, colframe=codeframe,
  boxrule=0.6pt, arc=2pt,
  left=8pt,right=8pt,top=6pt,bottom=6pt,
  fontupper=\ttfamily\small
}

% tip / warning / note boxes
\newtcolorbox{tip}{
  enhanced, breakable, colback=green!4, colframe=green!45!black,
  boxrule=0.8pt, arc=3pt, left=10pt,right=8pt,top=6pt,bottom=6pt,
  overlay unbroken and first={
    \node[anchor=north west] at ([xshift=-2pt,yshift=2pt]frame.north west){};}
}
\newtcolorbox{warn}{
  enhanced, breakable, colback=red!4, colframe=red!55!black,
  boxrule=0.8pt, arc=3pt, left=10pt,right=8pt,top=6pt,bottom=6pt
}
\newtcolorbox{note}{
  enhanced, breakable, colback=accent!5, colframe=accent,
  boxrule=0.8pt, arc=3pt, left=10pt,right=8pt,top=6pt,bottom=6pt
}

\setlist[itemize]{leftmargin=1.4em,itemsep=2pt,topsep=3pt}
\setlist[enumerate]{leftmargin=1.6em,itemsep=2pt,topsep=3pt}

\renewcommand{\arraystretch}{1.35}

%==================== TITLE BLOCK ====================
\begin{document}
\thispagestyle{empty}

\begin{tcolorbox}[enhanced,colback=poliblue,colframe=poliblue,
  arc=4pt,boxrule=0pt,left=18pt,right=18pt,top=16pt,bottom=16pt]
  {\color{white}\Huge\bfseries Politecnico di Milano}\\[4pt]
  {\color{polilight}\LARGE MSc Thesis Template --- Quick Guide}\\[10pt]
  {\color{white}\normalsize Everything you need to know, plus useful commands.}
\end{tcolorbox}

\vspace{4pt}
{\color{soft}\small This template is the official \emph{PoliMi3i} class
(School of Industrial and Information Engineering). You only edit
\code{Thesis.tex} and \code{Thesis\_bibliography.bib} --- never the files in
\code{Configuration\_Files/}.}

%==================== 1. STRUCTURE ====================
\section{Project structure}

\begin{tabular}{@{}>{\ttfamily}l l@{}}
\toprule
\normalfont\textbf{File / folder} & \textbf{What it is} \\
\midrule
Thesis.tex                & Main file --- \textbf{edit this} \\
Thesis\_bibliography.bib  & Your references (BibTeX database) \\
Configuration\_Files/PoliMi3i\_thesis.cls & Document class --- do not touch \\
Configuration\_Files/config.tex & Layout settings --- do not touch \\
Images/                   & Put your figures here \\
\bottomrule
\end{tabular}

%==================== 2. COMPILING ====================
\section{How to compile}

The bibliography uses BibTeX, so a single pass is \emph{not} enough. Run the
full sequence:

\begin{cmd}
pdflatex Thesis\\
bibtex   Thesis\\
pdflatex Thesis\\
pdflatex Thesis
\end{cmd}

\begin{note}
\marker{accent}{NOTE}\quad\textbf{On Overleaf} (this project was imported from it): just press
\textbf{Recompile}. Set \emph{Compiler} = \textbf{pdfLaTeX} and \emph{Main document}
= \code{Thesis.tex} in \emph{Menu}.
\end{note}

\begin{tip}
\marker{green!45!black}{TIP}\quad\textbf{Locally} you can use \code{latexmk -pdf Thesis.tex} to run the
whole sequence automatically, or \code{tectonic Thesis.tex} (resolves citations
and references in one command).
\end{tip}

%==================== 3. FILL FIRST ====================
\section{The first things to fill in}

\subsection{Title page --- \texttt{\textbackslash puttitle\{\}}}
\begin{cmd}
\textbackslash puttitle\{\\
\ \ title=My Thesis Title,\\
\ \ name=Name Surname,\\
\ \ course=Xxx Engineering - Ingegneria Xxx,\\
\ \ ID = 000000,\ \ \% numero di matricola\\
\ \ advisor= Prof. Name Surname,\\
\ \ coadvisor=\{Name Surname\},\ \ \% delete line if none\\
\ \ academicyear=\{2025-26\},\\
\}
\end{cmd}

\subsection{Abstracts (both required)}
PoliMi requires an abstract in \textbf{English} and one in \textbf{Italian}
(\emph{Abstract in lingua italiana}), each followed by 4--6 keywords.

\subsection{Your chapters}
Replace the demo chapters between \code{\textbackslash mainmatter} and the
bibliography with your own content.

%==================== 4. USEFUL COMMANDS ====================
\section{Useful commands cheat-sheet}

\begin{tabular}{@{}>{\ttfamily}p{0.50\textwidth} p{0.42\textwidth}@{}}
\toprule
\normalfont\textbf{Command} & \textbf{What it does} \\
\midrule
\textbackslash chapter\{Title\}              & New numbered chapter \\
\textbackslash chapter*\{Title\}             & Unnumbered chapter \\
\textbackslash section\{...\} / \textbackslash subsection\{...\} & Sections (use \texttt{*} for no number) \\
\textbackslash input\{Chapters/ch1.tex\}     & Pull in a chapter from another file \\
\textbackslash label\{ch:intro\}             & Set a reference anchor \\
\textbackslash ref\{ch:intro\}               & Reference it (auto via cleveref) \\
\textbackslash eqref\{eq:maxwell\}           & Reference an equation \\
\textbackslash cite\{key\}                    & Cite a .bib entry \\
\textbackslash includegraphics[width=...]\{f\} & Insert an image from Images/ \\
\textbackslash frontmatter / \textbackslash mainmatter & Roman (i,ii) then arabic (1,2) pages \\
\bottomrule
\end{tabular}

\subsection{Label naming convention}
Stay consistent so cross-references read well:

\begin{center}
\begin{tabular}{@{}>{\ttfamily}l l@{}}
\toprule
\normalfont\textbf{Prefix} & \textbf{For} \\
\midrule
ch:    & chapters \\
sec:   & sections \\
fig:   & figures \\
table: & tables \\
eq:    & equations \\
alg:   & algorithms \\
\bottomrule
\end{tabular}
\end{center}

%==================== 5. BUILDING BLOCKS ====================
\section{Building blocks (already demoed in the template)}

\subsection{Figure}
\begin{cmd}
\textbackslash begin\{figure\}[H]\\
\ \ \textbackslash centering\\
\ \ \textbackslash includegraphics[width=0.6\textbackslash textwidth]\{myplot.png\}\\
\ \ \textbackslash caption\{My caption.\}\\
\ \ \textbackslash label\{fig:myplot\}\\
\textbackslash end\{figure\}
\end{cmd}

\subsection{Table}
\begin{cmd}
\textbackslash begin\{table\}[H]\\
\ \ \textbackslash centering\\
\ \ \textbackslash begin\{tabular\}\{|p\{3em\} c c c|\}\\
\ \ \textbackslash hline\\
\ \ \textbackslash rowcolor\{bluepoli!40\}\ \% PoliMi blue header\\
\ \ \ \& \textbackslash textbf\{col 1\} \& \textbackslash textbf\{col 2\} \& \textbackslash textbf\{col 3\}\textbackslash\textbackslash\\
\ \ \textbackslash hline\textbackslash hline\\
\ \ \textbackslash textbf\{row 1\} \& 1 \& 2 \& 3 \textbackslash\textbackslash\\
\ \ \textbackslash hline\\
\ \ \textbackslash end\{tabular\}\\
\ \ \textbackslash caption\{My caption.\} \textbackslash label\{table:ex\}\\
\textbackslash end\{table\}
\end{cmd}

\subsection{Equation}
\begin{cmd}
\textbackslash begin\{equation\}\\
\ \ \textbackslash label\{eq:energy\}\\
\ \ E = m c\textasciicircum 2\\
\textbackslash end\{equation\}
\end{cmd}

%==================== 6. BIBLIOGRAPHY ====================
\section{Bibliography}

Add entries to \code{Thesis\_bibliography.bib}, then cite with \code{\textbackslash cite\{key\}}:

\begin{cmd}
@article\{einstein1905,\\
\ \ author  = \{A. Einstein\},\\
\ \ title   = \{On the Electrodynamics of Moving Bodies\},\\
\ \ journal = \{Annalen der Physik\},\\
\ \ year    = \{1905\}\\
\}
\end{cmd}

Style is \code{abbrvnat} (numbered, \code{[1]}). Change \code{\textbackslash bibliographystyle}
in \code{Thesis.tex} if your field prefers another.

%==================== 7. TIPS ====================
\section{Good-to-know tips}

\begin{itemize}
  \item \textbf{Split long theses:} one \code{.tex} file per chapter, included via
        \code{\textbackslash input\{\}}. Easier to manage and compiles faster.
  \item \textbf{PoliMi blue} is predefined as \code{bluepoli} --- use it for table headers.
  \item \textbf{Cell padding:} the \code{\textbackslash T} and \code{\textbackslash B} macros add
        top/bottom space inside table cells.
  \item \textbf{Fixed structure:} the order (Abstract EN+IT $\rightarrow$ ToC $\rightarrow$
        chapters $\rightarrow$ Bibliography $\rightarrow$ Appendices $\rightarrow$ Lists
        $\rightarrow$ Acknowledgements) is set by PoliMi rules --- keep it.
  \item \textbf{Image formats:} \code{.pdf}, \code{.png}, \code{.jpg}, \code{.eps} all work
        with pdfLaTeX. Vector (\code{.pdf}/\code{.eps}) looks sharpest.
\end{itemize}

\begin{warn}
\marker{red!55!black}{WARNING}\quad\textbf{Don't edit} \code{PoliMi3i\_thesis.cls} or
\code{config.tex}. They control the official formatting; changing them risks a
non-compliant thesis.
\end{warn}

\begin{tip}
\marker{green!45!black}{TIP}\quad The whole \code{Thesis.tex} is a \textbf{live tutorial} --- every section
demonstrates the feature it describes. Compile it once as-is, then replace the demo
text piece by piece.
\end{tip}

\vfill
\begin{center}
{\color{soft}\small Generated for your \texttt{Tex\_tesi\_vibra} project \textbullet\ Keep this handy while writing.}
\end{center}

\end{document}
